\documentclass[11pt,letterpaper]{article}

% ------------------------------------------------------------
% Paquetes
% ------------------------------------------------------------
\usepackage[utf8]{inputenc}
\usepackage[T1]{fontenc}
\usepackage{lmodern}
\usepackage[spanish,es-tabla,es-nodecimaldot]{babel}
\usepackage[margin=2.5cm]{geometry}
\usepackage{amsmath,amssymb}
\usepackage{graphicx}
\usepackage{booktabs}
\usepackage{tabularx}
\usepackage{float}
\usepackage{caption}
\usepackage{xcolor}
\usepackage{tikz}
\usetikzlibrary{arrows.meta,positioning,fit,babel}
\usepackage{url}
\usepackage[hidelinks]{hyperref}

\captionsetup{font=small,labelfont=rm}
\setlength{\parskip}{0.3em}
\setlength{\emergencystretch}{2em}

% Enlace al repositorio (cambiar por la URL real)
\newcommand{\repourl}{https://github.com/SebasDeb/practica4-uart-fpga-esp32-matlab}

% ------------------------------------------------------------
\begin{document}

% ============================================================
% Portada
% ============================================================
\begin{titlepage}
\centering
\vspace*{1.5cm}

{\large Universidad de las Américas Puebla (UDLAP)\par}
\vspace{0.3cm}
{\large Laboratorio de Sistemas Embebidos\par}

\vspace{2.5cm}

{\Large\bfseries Reporte de Práctica 4\par}
\vspace{0.6cm}
{\LARGE Comunicación bidireccional UART entre la FPGA DE10-Lite y MATLAB
mediante un ESP32\par}
\vspace{0.6cm}
{\Large Cálculo de $A^{n/2}$ en formato Q80.20\par}

\vspace{2.5cm}

{\large
Abril Morales Huerta\\[0.2em]
Anellisse Herrera Maldonado\\[0.2em]
José Manuel Daniel\par}

\vfill

{\small Repositorio del proyecto:\\
\url{\repourl}\par}

\vspace{1cm}
{\large 30 Octubre 2026\par}
\end{titlepage}

% ============================================================
\section{Resumen de la práctica}
% ============================================================

En esta práctica diseñamos e implementamos un sistema de comunicación serial
bidireccional entre una computadora y la tarjeta FPGA DE10-Lite. El usuario
introduce en MATLAB una base $A$, de 0 a 255, y un exponente $n$, de 0 a 15.
MATLAB envía ambos valores como dos bytes. Un ESP32, programado como puente
transparente USB--UART, los reenvía a la FPGA. La FPGA calcula $A^{n/2}$ en
formato de punto fijo Q80.20 y regresa el resultado de 100 bits dividido en
13 bytes. Finalmente, MATLAB reconstruye el valor, lo convierte a decimal y
lo compara con una referencia calculada por software.

La comunicación utiliza 9600 baudios con formato 8N1. A partir del reloj de
50~MHz, la FPGA cuenta $50\,000\,000/9600 \approx 5208$ ciclos por bit, lo que
produce un error de velocidad de 0.0064\%. El ESP32 y la FPGA trabajan con
niveles de 3.3~V, por lo que se conectan directamente: el GPIO17 del ESP32
con la entrada RX de la FPGA, el GPIO16 con la salida TX y ambas tierras
entre sí. La Figura~\ref{fig:arquitectura} muestra la arquitectura del
sistema.

\begin{figure}[H]
\centering
\begin{tikzpicture}[
  >=Stealth,
  font=\scriptsize,
  blk/.style={draw,thick,rounded corners=2pt,align=center,minimum height=0.9cm,minimum width=1.9cm},
  grp/.style={draw,dashed,rounded corners=4pt,inner sep=6pt}
]
  \node[blk] (pc) at (0,0) {MATLAB\\(PC)};
  \node[blk] (esp) at (3,0) {ESP32\\UART0 / UART2};
  \node[blk] (rx) at (6.6,1.1) {\texttt{uart\_rx}};
  \node[blk] (tx) at (6.6,-1.1) {\texttt{uart\_tx}};
  \node[blk] (fsm) at (9.4,0) {Control\\\texttt{top\_uart}\\\texttt{\_de10lite}};
  \node[blk] (pot) at (12.6,0) {\texttt{potencia}\\\texttt{\_q8020}};
  \node[grp,fit=(rx)(tx)(fsm)(pot),label={[font=\scriptsize]above:FPGA DE10-Lite (50 MHz)}] {};
  \draw[<->,thick] (pc) -- node[above] {USB} (esp);
  \draw[->,thick] (esp.north) |- node[pos=0.75,above] {GPIO17 (TX2)} (rx.west);
  \draw[<-,thick] (esp.south) |- node[pos=0.75,below] {GPIO16 (RX2)} (tx.west);
  \draw[->] (rx.east) -| node[pos=0.25,above] {byte} (fsm.north);
  \draw[->] (fsm.south) |- node[pos=0.75,below] {byte} (tx.east);
  \draw[->] ([yshift=2mm]fsm.east) -- node[above] {$A$, $n$} ([yshift=2mm]pot.west);
  \draw[<-] ([yshift=-2mm]fsm.east) -- node[below] {100 bits} ([yshift=-2mm]pot.west);
  \node[font=\scriptsize] at (3,-2) {GND común entre ESP32 y FPGA};
\end{tikzpicture}
\caption{Arquitectura del sistema. Elaboración propia.}
\label{fig:arquitectura}
\end{figure}

En la FPGA organizamos el diseño en cuatro módulos Verilog. El receptor
\texttt{uart\_rx} sincroniza la línea con dos flip-flops, confirma el bit de
inicio a la mitad de su duración, muestrea los ocho bits de datos y descarta
las tramas con un bit de paro inválido. El transmisor \texttt{uart\_tx} genera
la trama 8N1 de cada byte. El módulo \texttt{potencia\_q8020} descompone el
exponente como $n = 2k + r$ y calcula
\begin{equation}
A^{n/2} = A^{k}\left(\sqrt{A}\right)^{r},
\end{equation}
mediante $k$ multiplicaciones sucesivas y una raíz cuadrada entera de
$A \cdot 2^{40}$, que entrega $\sqrt{A}$ con 20 bits fraccionarios. El
cálculo tarda como máximo 9 ciclos de reloj. Por último, el módulo
\texttt{top\_uart\_de10lite} coordina la secuencia con una máquina de
estados: recibe $A$ y $n$, inicia el cálculo y envía los 13 bytes del
resultado, comenzando por el menos significativo.

Verificamos el módulo de cálculo en Active-HDL con un banco de pruebas
autoverificable. Los cinco casos aplicados produjeron el valor esperado. En
la tarjeta probamos la comunicación completa desde MATLAB. La
Tabla~\ref{tab:resultados} muestra algunos resultados representativos.

\begin{table}[H]
\centering
\caption{Resultados representativos del cálculo de $A^{n/2}$.}
\label{tab:resultados}
\small
\begin{tabular}{@{}rrrrr@{}}
\toprule
$A$ & $n$ & \textbf{Valor exacto} & \textbf{Resultado Q80.20} & \textbf{Error relativo} \\
\midrule
4   & 3  & 8           & 8           & 0 \\
100 & 4  & 10\,000     & 10\,000     & 0 \\
2   & 1  & 1.414213562 & 1.414213181 & $2.7\times10^{-7}$ \\
10  & 5  & 316.2277660 & 316.2277222 & $1.4\times10^{-7}$ \\
255 & 15 & $1.119570259\times10^{18}$ & $1.119570197\times10^{18}$ & $5.6\times10^{-8}$ \\
\bottomrule
\end{tabular}
\end{table}

Con exponentes pares, el resultado es exacto. Con exponentes impares, la
raíz se trunca a 20 bits fraccionarios, por lo que el error relativo es
menor que $1/(\sqrt{2}\cdot2^{20}) \approx 6.7\times10^{-7}$. Cada operación
completa tarda aproximadamente 15.6~ms, casi todo debido a la transmisión
serial.

La práctica nos permitió comprobar la importancia de que todos los
elementos usen el mismo baud rate, la utilidad del ESP32 como sustituto de
un convertidor USB--UART dedicado y las ventajas de MATLAB para generar
casos de prueba y validar resultados. El código fuente, las simulaciones, la
documentación técnica completa, las capturas y el video de funcionamiento se
encuentran en el repositorio del proyecto:

\begin{center}
\url{\repourl}
\end{center}

% ============================================================
\section{Cuestionario}
% ============================================================

\begin{enumerate}
  \item \textbf{¿Qué diferencia existe entre UART y RS232?}

  UART es el bloque que implementa la comunicación serial asíncrona: define la
  trama, con sus bits de inicio, datos, paridad y paro, y trabaja con los
  niveles lógicos del circuito, como 3.3~V en nuestro caso.

  RS232 es un estándar eléctrico y de interfaz. Define niveles de voltaje,
  conectores y señales de control. Utiliza voltajes positivos y negativos,
  típicamente entre $\pm3$ y $\pm15$~V, con lógica invertida: el 1 lógico
  corresponde a un voltaje negativo \cite{ti232}. Una UART puede usarse con
  RS232 agregando un transceptor que convierta los niveles.

  \item \textbf{¿Por qué no debe conectarse directamente una señal RS232
  clásica a un pin GPIO de la FPGA?}

  Los pines de la DE10-Lite que utilizamos operan con 3.3~V \cite{de10}. Una
  señal RS232 puede alcanzar voltajes de varios volts, positivos y negativos,
  que superan los límites del pin y pueden dañar la FPGA. Además, su lógica
  invertida no corresponde a la que espera nuestro receptor. Por ello, se
  requiere un convertidor de niveles. En nuestro sistema, el ESP32 y la FPGA
  trabajan con los mismos niveles, por lo que no fue necesario.

  \item \textbf{¿Qué función cumple el baud rate en una comunicación serial?}

  Define el número de símbolos transmitidos por segundo y, por tanto, la
  duración de cada bit. Como no existe un reloj compartido, el receptor usa
  esa duración para saber en qué momento muestrear cada bit después del bit
  de inicio. Si las velocidades no coinciden, el error se acumula a lo largo
  de la trama y el receptor puede leer un bit equivocado \cite{brown}.

  En nuestro diseño, 9600 baudios corresponden a 5208 ciclos de 50~MHz por
  bit, con un error de 0.0064\%.

  \item \textbf{¿Por qué la conexión entre TX y RX debe realizarse cruzada?}

  TX es una salida que transmite y RX es una entrada que recibe. Al conectar
  el TX de un dispositivo con el RX del otro, cada transmisor tiene un
  receptor que lo escucha. Si se conectara TX con TX, dos salidas manejarían
  la misma línea y ningún dispositivo recibiría los datos. En nuestro sistema,
  el GPIO17 del ESP32 se conecta a la entrada RX de la FPGA y el GPIO16 a su
  salida TX.

  \item \textbf{¿Qué ventajas ofrece MATLAB para la prueba de sistemas
  embebidos con comunicación serial?}

  El objeto \texttt{serialport} permite abrir el puerto, configurar el tiempo
  límite, limpiar el búfer, enviar y leer bytes con pocas instrucciones
  \cite{matlab}. Además, MATLAB permite automatizar casos de prueba, validar
  las entradas, procesar los datos recibidos y compararlos con valores de
  referencia.

  En esta práctica, la aritmética simbólica nos permitió reconstruir un
  resultado de 100 bits sin pérdida de precisión.

  \item \textbf{¿Qué función cumple el ESP32 dentro del sistema de
  comunicación propuesto?}

  Funciona como un puente transparente USB--UART \cite{esp32}. Recibe por USB
  los bytes que envía MATLAB y los retransmite a la FPGA por su UART2. En
  sentido contrario, recibe la respuesta de la FPGA y la reenvía a la
  computadora. Así, MATLAB se comunica con la FPGA como si estuviera conectada
  a un puerto serial, sin requerir un convertidor USB--UART dedicado.
\end{enumerate}

% ============================================================
\begin{thebibliography}{9}
% ============================================================

\bibitem{mano}
M. M. Mano, \textit{Diseño digital}, versión en español proporcionada para
consulta.

\bibitem{de10}
Terasic, \textit{DE10-Lite User Manual}, Jun. 2020.

\bibitem{esp32}
Espressif Systems, \textit{ESP32 Technical Reference Manual}, capítulo UART
Controller. [Online]. Available:
\url{https://www.espressif.com/sites/default/files/documentation/esp32_technical_reference_manual_en.pdf}

\bibitem{matlab}
MathWorks, \textit{serialport: Connection to serial port}, documentación de
MATLAB. [Online]. Available:
\url{https://www.mathworks.com/help/matlab/ref/serialport.html}

\bibitem{brown}
S. Brown y Z. Vranesic, \textit{Fundamentals of Digital Logic with Verilog
Design}, 3rd ed. New York, NY, USA: McGraw-Hill, 2014.

\bibitem{ti232}
Texas Instruments, \textit{Interface Circuits for TIA/EIA-232-F}, Application
Report SLLA037.

\end{thebibliography}

\end{document}
